\begin{figure}[H]
    \centering
    \begin{tikzpicture}[scale=1.35, >=Latex]
        \draw[->] (-2.8,0) -- (2.8,0) node[right] {$x$};
        \draw[->] (0,0.15) -- (0,2.05) node[above] {$f(x)$};

        \fill[blue!10] (-0.75,0.75) rectangle (2.45,1.25);
        \draw[blue!55!black, dashed] (-2.45,0.75) -- (2.45,0.75)
            node[right] {$3/4$};
        \draw[blue!55!black, dashed] (-2.45,1.25) -- (2.45,1.25)
            node[right] {$5/4$};
        \node[blue!55!black, anchor=west] at (0.25,1.03) {$V=(3/4,5/4)$};

        \draw[orange!85!black, very thick] (-0.75,1.75) -- (0,1);
        \draw[orange!85!black, very thick] (0.02,0.48) -- (0.95,-0.45);

        \fill[orange!85!black] (0,1) circle (2pt);
        \draw[orange!85!black, fill=white, thick] (0,0.5) circle (2pt);

        \draw[dashed] (0,0.15) -- (0,1.95);
        \node[below left] at (0,0) {$0$};

        \draw[red!70!black, ultra thick] (-0.25,0) -- (0,0);
        \draw[red!70!black, fill=white, thick] (-0.25,0) circle (1.6pt);
        \fill[red!70!black] (0,0) circle (1.8pt);
        \node[red!70!black, anchor=north] at (-0.72,-0.18) {$f^{-1}(V)=(-1/4,0]$};
    \end{tikzpicture}
    \caption{For the jump at $0$, the open interval $V=(3/4,5/4)$ contains $f(0)=1$ but not where the graph continues near $1/2$. Its preimage is $(-1/4,0]$, not an open subset of $\mathbb{R}$, matching the calculus intuition that the function should have discontinuity at 0.}
    \label{fig:continuous-map-preimage-calculus}
\end{figure}
